\section{构建完整的 CNN}

\subsection{经典 CNN 架构模式}

\begin{figure}[H]
\centering
\includegraphics[width=0.95\textwidth]{slides-images/slide-060.jpg}
\caption{典型 CNN 架构：卷积层、激活函数和池化层交替出现，最后是全连接层\protect\footnotemark}
\end{figure}
\footnotetext{Slides 第61页。}

经典的 CNN 遵循以下模式：

$$\text{[Conv} \rightarrow \text{ReLU} \rightarrow \text{Pool]}^N \rightarrow \text{[FC} \rightarrow \text{ReLU]}^M \rightarrow \text{FC (输出)}$$

\begin{itemize}
    \item 随着网络深度增加，空间尺寸逐渐\textbf{缩小}（通过池化或带步幅的卷积）
    \item 通道数逐渐\textbf{增加}（每次池化后通常将通道数翻倍）
    \item 直觉：空间尺寸缩小但通道数增加，总信息量大致保持不变
\end{itemize}

\subsection{详细示例：一个简单的 CNN}

\begin{figure}[H]
\centering
\includegraphics[width=0.95\textwidth]{slides-images/slide-061.jpg}
\caption{完整 CNN 示例：$224 \times 224$ 输入经过多层卷积和池化，最终输出 1000 类分数\protect\footnotemark}
\end{figure}
\footnotetext{Slides 第62页。}

以输入尺寸 $3 \times 224 \times 224$ 的 CNN 为例，追踪各层的张量形状变化：

\begin{table}[H]
\centering
\small
\begin{tabular}{lccc}
\toprule
\textbf{层} & \textbf{输出形状} & \textbf{参数量} & \textbf{说明} \\
\midrule
输入 & $3 \times 224 \times 224$ & -- & RGB 图像 \\
Conv1 ($7 \times 7$, S=2) & $64 \times 112 \times 112$ & 9,472 & 下采样 2x \\
Pool ($2 \times 2$, S=2) & $64 \times 56 \times 56$ & 0 & 下采样 2x \\
Conv2 ($3 \times 3$) & $192 \times 56 \times 56$ & 110,784 & 增加通道 \\
Pool ($2 \times 2$, S=2) & $192 \times 28 \times 28$ & 0 & 下采样 2x \\
Conv3 ($3 \times 3$) & $384 \times 28 \times 28$ & 663,936 & 增加通道 \\
\bottomrule
\end{tabular}
\caption{CNN 各层张量形状变化示例}
\end{table}

\subsection{处理不同尺寸的输入}

\begin{warningbox}{输入图像必须大小一致}
在当前框架下，同一个 batch 中的所有图像必须具有\textbf{相同的空间尺寸}。常见的处理方法：
\begin{itemize}
    \item 将所有图像 resize 到固定大小（最常见）
    \item 用零填充到统一大小
    \item 宽高比分桶（Aspect Ratio Bucketing）——同一 batch 内使用相同宽高比，不同 batch 可以不同
\end{itemize}
\end{warningbox}

\subsection{本章小结}

构建一个完整的 CNN 就是在计算图中组合卷积层、激活函数、池化层和全连接层。核心设计原则是：空间尺寸逐层缩小，通道数逐层增加，形成从像素到语义的特征金字塔。输出端用全连接层（或全局平均池化 + 线性层）将特征映射到分类分数。

%% ========== Part 8: AlexNet ==========

